% Generated draft: requires mathematical and source review.
\documentclass[11pt]{article}
\usepackage[margin=25mm]{geometry}
\usepackage{fontspec}
\setmainfont{Latin Modern Roman}
\usepackage{amsmath,amssymb,mathtools,unicode-math}
\setmathfont{Latin Modern Math}
\usepackage{enumitem,needspace,xurl,hyperref}
\setlength{\parindent}{0pt}
\setlength{\parskip}{5pt}
\setlength{\emergencystretch}{3em}
\newcommand{\N}{\mathbb{N}}
\newcommand{\Z}{\mathbb{Z}}
\newcommand{\Q}{\mathbb{Q}}
\newcommand{\R}{\mathbb{R}}
\newcommand{\C}{\mathbb{C}}
\newcommand{\F}{\mathbb{F}}
\newcommand{\A}{\mathbb{A}}
\newcommand{\PP}{\mathbb{P}}
\newcommand{\E}{\mathbb{E}}
\newcommand{\eps}{\varepsilon}
\newcommand{\dd}{\,\mathrm{d}}
\begin{document}
\section*{UnsolvedMath ID: 30003425}
\textbf{Source problem code:} catalogue entry
\subsection{Definitions and mathematical statement}
For a specified option-pricing convention and maturity \(T\), the implied volatility at strike \(K\) is the volatility parameter reproducing the option price, whenever that parameter exists and is unique. The volatility smile is the function \(K\mapsto\sigma_{\mathrm{imp}}(K,T)\). The source does not specify the pricing convention, admissible strikes, maturities, or conditions ensuring that implied volatility is defined.

To describe the requested inference without imposing additional hypotheses, let \(\mathcal K_T\) denote an unspecified admissible strike domain and let \(\mathcal O_T\subseteq\mathcal K_T\) denote its observable strikes. Write \(P\) for a candidate stochastic model, \(\sigma_P(K,T)\) for its implied volatility where defined, and \(\Theta(P)\) for the intended characteristic of the stochasticity of volatility. The model class and the meaning of \(\Theta\) are not supplied. The observations would consist of market estimates of \(\sigma_P(K,T)\) for \(K\in\mathcal O_T\); extrapolation concerns values for \(K\in\mathcal K_T\setminus\mathcal O_T\).

\textbf{Source question.} Can the volatility smile be extrapolated from observable strikes robustly enough to infer accurately the stochasticity of volatility from market data?

In the notation above, the requested reconstruction has the schematic form
\[
 \bigl(\sigma_P(K,T)\bigr)_{K\in\mathcal O_T}
 \longmapsto
 \left(\bigl(\sigma_P(K,T)\bigr)_{K\in\mathcal K_T\setminus\mathcal O_T},\,\Theta(P)\right).
\]
This display describes the intended inference task; it does not assert that such a reconstruction exists or is unique. The source supplies no observation-error model, admissible model class, reconstruction accuracy, stability criterion, or quantifiers over models and market data. In particular, it does not define whether stochasticity means randomness itself, a volatility-of-volatility parameter, path regularity, or a law of a volatility process. These unresolved choices prevent a unique fully quantified mathematical statement. The notation at maturity \(T\) does not prescribe whether data from one maturity or several maturities are intended.
\subsection{Short English statement}
The question asks whether observed portions of a volatility smile permit reliable extrapolation and accurate inference about the randomness of volatility. The source leaves the model class, inference target, data errors, and meanings of robustness and accuracy unspecified, so this remains a research question requiring further hypotheses and human mathematical review.
\subsection{Sources}
\begin{itemize}
\item Supplied local UnsolvedMath registry record, permanent ID \(30003425\), exact source problem code \texttt{catalogue entry}, titled ``Volatility Smile Extrapolation and Latent Stochasticity,'' in the Oberwolfach Reports Open Problems collection. No individual proposer is identified; the supplied proposed year is 2017.
\item Supplied report citation: \emph{Mathematics of Quantitative Finance (2018)}, Oberwolfach Reports, DOI 10.4171/owr/2017/13, pp.\ 683--769. \url{https://doi.org/10.4171/owr/2017/13}. The record describes Oberwolfach Reports as documenting mathematical workshops and their open-problem discussions. The citation's 2018 label and its 2017 DOI component are retained as supplied; their relationship has not been verified.
\item Supplied literature reference: Gatheral--Jaisson--Rosenbaum, ``rough volatility evidence.'' \url{https://arxiv.org/abs/1702.02777}. This is the supplied descriptive label; the bibliographic identity and its support for the assessment have not been independently verified.
\item Supplied literature reference: Guyon--Lekeufack, ``volatility models in practice.'' \url{https://arxiv.org/abs/2401.03345}. This is the supplied descriptive label; the bibliographic identity and its support for the assessment have not been independently verified.
\item Literature assessment dated exactly \texttt{2026-08-21}: the supplied curation labels the problem open and states that rough and path-dependent models improve smile fitting and extrapolation diagnostics, while reporting no model-independent result showing that visible strikes robustly identify latent stochasticity or determine accurate extrapolation. This is supplied dated literature triage derived from the OWR curation, not an independently verified literature assessment or mathematical fact; its time-sensitive status and cited claims require independent verification.
\item The registry labels the statement \texttt{corrected\_verified} and reports that mathematical content was preserved while notation, typography, and sentence structure were repaired for publication. This is supplied editorial verification metadata, not independent verification of the formulation or its open status.
\end{itemize}
\end{document}
